\bibitem{clavel1996reflection}
M. Clavel and J. Meseguer. Reflection and Strategies in Rewriting Logic. \emph{Electronic Notes in Theoretical Computer Science} 4:126--148, 1996. doi:10.1016/S1571-0661(04)00037-4.

\bibitem{rosu2010k}
G. Ro\c{s}u and T. F. \c{S}erb\u{a}nu\c{t}\u{a}. An overview of the K semantic framework. \emph{Journal of Logic and Algebraic Programming} 79(6):397--434, 2010. doi:10.1016/j.jlap.2010.03.012.

\bibitem{demoura2021lean}
L. de Moura and S. Ullrich. The Lean 4 Theorem Prover and Programming Language. In \emph{Automated Deduction -- CADE 28}, LNCS 12699, pages 625--635, 2021. doi:10.1007/978-3-030-79876-5\_37.

\bibitem{necula1997pcc}
G. C. Necula. Proof-carrying code. In \emph{Proceedings of POPL 1997}, pages 106--119, 1997. doi:10.1145/263699.263712.

\bibitem{clarke2000cegar}
E. M. Clarke, O. Grumberg, S. Jha, Y. Lu, and H. Veith. Counterexample-Guided Abstraction Refinement. In \emph{Computer Aided Verification (CAV 2000)}, LNCS 1855, pages 154--169, 2000. doi:10.1007/10722167\_15.

\bibitem{blackwell1953}
D. Blackwell. Equivalent Comparisons of Experiments. \emph{Annals of Mathematical Statistics} 24(2):265--272, 1953. doi:10.1214/aoms/1177729032.

\bibitem{hay2014computations}
N. Hay, S. Russell, D. Tolpin, and S. E. Shimony. Selecting Computations: Theory and Applications. arXiv:1408.2048, 2014.

\bibitem{ramdas2023}
A. Ramdas, P. Gr\"unwald, V. Vovk, and G. Shafer. Game-Theoretic Statistics and Safe Anytime-Valid Inference. \emph{Statistical Science} 38(4):576--597, 2023. doi:10.1214/23-STS894.

\bibitem{dekleer1986}
J. de Kleer. An assumption-based TMS. \emph{Artificial Intelligence} 28(2):127--162, 1986. doi:10.1016/0004-3702(86)90080-9.

\bibitem{dung1995}
P. M. Dung. On the acceptability of arguments and its fundamental role in nonmonotonic reasoning, logic programming and n-person games. \emph{Artificial Intelligence} 77(2):321--357, 1995. doi:10.1016/0004-3702(94)00041-X.

\bibitem{schmidhuber2007godel}
J. Schmidhuber. G\"odel Machines: Fully Self-referential Optimal Universal Self-improvers. In B. Goertzel and C. Pennachin, editors, \emph{Artificial General Intelligence}, pages 199--226. Springer, 2007. doi:10.1007/978-3-540-68677-4\_7.

\bibitem{fallenstein2014selfref}
B. Fallenstein and N. Soares. Problems of Self-reference in Self-improving Space-Time Embedded Intelligence. In \emph{Artificial General Intelligence 2014}, LNCS 8598, pages 21--32, 2014. doi:10.1007/978-3-319-09274-4\_3.

\bibitem{boiko2023coscientist}
D. A. Boiko, R. MacKnight, B. Kline, and G. Gomes. Autonomous chemical research with large language models. \emph{Nature} 624:570--578, 2023. doi:10.1038/s41586-023-06792-0.

\bibitem{bran2024chemcrow}
A. M. Bran, S. Cox, O. Schilter, C. Baldassari, A. D. White, and P. Schwaller. Augmenting large language models with chemistry tools. \emph{Nature Machine Intelligence} 6:525--535, 2024. doi:10.1038/s42256-024-00832-8.

\bibitem{gottweis2026coscientist}
J. Gottweis et al. Accelerating scientific discovery with Co-Scientist. \emph{Nature} 655:487--496, 2026. doi:10.1038/s41586-026-10644-y.

\bibitem{ghareeb2026multiagent}
A. Ghareeb et al. A multi-agent system for automating scientific discovery. \emph{Nature} 655:497--505, 2026. doi:10.1038/s41586-026-10652-y.
